\documentclass[10pt]{amsart}
\usepackage{amssymb,amsmath,amsthm,amsfonts,amsopn,url}
\usepackage{amscd,amssymb,amsopn,amsmath,amsthm,graphics,amsfonts,enumerate,verbatim,calc}
\usepackage[all]{xy}
\theoremstyle{plain}
\newtheorem{thm}{Theorem}[section]
\newtheorem*{mt*}{Theorem (Huneke-Lyubeznik)}
\newtheorem*{cj*}{Conjecture}
\newtheorem*{nt*}{Notations}
\newtheorem*{stfac}{Base change formulas for depth}
\newtheorem{prop}[thm]{Proposition}
\newtheorem{lemma}[thm]{Lemma}
\newtheorem{cor}{Corollary}
\newtheorem{rem}{Remark}
\newtheorem{example}{Example}[section]
\theoremstyle{definition}
\newtheorem*{definition}{Definition (The graph $\Theta_R$)}
\newcommand{\ideal}[1]{\mathfrak{#1}}
\newcommand{\M}{\ideal{M}}
\newcommand{\m}{\ideal{m}}
\newcommand{\n}{\ideal{n}}
%\newcommand{\o}{\ideal{o}}
\newcommand{\p}{\ideal{p}}
\newcommand{\q}{\ideal{q}}
\newcommand{\I}{\ideal{I}}
\newcommand{\J}{\ideal{J}}
\newcommand{\ffunc}[1]{\mathrm{#1}}
\newcommand{\func}[1]{\mathrm{#1} \,}
\newcommand{\depth}{\func{depth}}
\newcommand{\depthsb}{\ffunc{depth}}
\newcommand{\height}{\func{ht}}
\newcommand{\grade}{\func{grade}}
\newcommand{\ara}{\func{ara}}
\newcommand{\gldim}{\func{gl} \dim}
\newcommand{\Tor}{\func{Tor}}
\newcommand{\Hom}{\func{Hom}}
\newcommand{\Ass}{\func{Ass}}
\newcommand{\im}{\func{im}}
\newcommand{\V}{\func{V}}
\newcommand{\rad}{\func{rad}}
\newcommand{\Supp}{\func{Supp}}
\newcommand{\Spec}{\func{Spec}}
\newcommand{\Ker}{\func{Ker}}
\newcommand{\Coker}{\func{Coker}}
\newcommand{\pdp}[1]{\ffunc{ph. depth}_{#1}}
\newcommand{\arrow}[1]{\stackrel{#1}{\rightarrow}}
\newcommand{\PP}{{\mathbb P}}
\newcommand{\QQ}{{\mathbb Q}}
\newcommand{\NN}{{\mathbb N}}
\newcommand{\ZZ}{{\mathbb Z}}
\newcommand{\CC}{{\mathbb C}}


\title[]{Finiteness of local cohomology modules over arbitrary regular local rings}

\author[]{Rajsekhar Bhattacharyya}

\address{Dinabandhu Andrews College, Garia, Kolkata 700084, India}

\email{rbhattacharyya@gmail.com}

\thanks{}

\keywords{Local Cohomology}

\subjclass[2010]{13D45}
\begin{document}

\begin{abstract}
In this paper, we consider arbitrary regular local ring $S$ and an ideal $J\subset S$ whose minimal primes are cohomologically complete intersection. Using the Mayer-Vietoris spectral sequence, in Section 2, we construct examples of local cohomology modules $H^i_J (S)$, whose set of associated primes is finite. As a special case we show that the set of associated primes $H^i_J (S)$ is finite when minimal primes of $J$ are part of a specific regular system of parameters and satisfy a `Dropping Condition'. 

In Section 3, we prove Theorem 2.3 along with its asssociated corollary in a more general setting, without the use of `Dropping Condition'. 

These examples of local cohomology modules include the cases whose associative primes also can contain prime number $p$.
\end{abstract}

\maketitle
\section{introduction}

Let $R$ be a Noetherian ring and $M$ be a module over it. For an ideal $I\subset R$ and for some integer $i\geq 0$, consider local cohomology module $H^i_I (M)$ with support in the ideal $I$. In the fourth problem of \cite{Hu}, it is asked that whether local cohomology modules of Noetherian rings have finitely many associated prime ideals. There are examples given in \cite{Si}, \cite{Ka}, and \cite{SS}, which show that the set of associated primes of $i$-th local cohomology module $H^i_I (R)$ of a Noetherian ring $R$ can be infinite. However, there are several important cases where we have the finiteness of the set of associated primes of local cohomology modules and we list them below: (1) Regular rings of prime characteristic \cite{HS}, (2) regular local and affine rings of characteristic zero \cite{Ly1}, (3) unramified regular local rings of mixed characteristic \cite{Ly2} and (4) smooth algebra over $\ZZ$ \cite{BBLSZ}. On the other direction, the result is also true for every low dimensional regular local ring upto dimension five \cite{Boix}. These results support Lyubeznik conjecture, (\cite{Ly1}, Remark 3.7): 

\textbf{Conjecture:} Let $R$ be a regular ring and $I\subset R$ be its ideal, then for each $i\geq 0$, $i$-th local cohomology module $H^i_I (R)$ has finitely many associated prime ideals.

Recently, in \cite{Ma} it is shown that in ramified regular local rings, there exist examples of local cohomology modules with ifinitely many associated primes. To investigate the cause in more detail, in this paper, we find a `Dropping Condition' (see definition inside) so that if the minimal prime ideals of the support ideal of a local cohomology module satisfy it, then its set of associated primes is finite. Moreover this result is true for any regular local ring. 

At first, in Section 2 (The Main Result I), we consider arbitrary regular local ring $S$ and an ideal $J\subset S$ whose minimal primes are cohomologically complete intersection \cite{HeSc}. Using the Mayer-Vietoris spectral sequence, we construct examples of local cohomology modules $H^i_J (S)$, whose set of associated primes is finite. For detail, see, Theorem 2.3. As a special case we show that the set of associated primes $H^i_J (S)$ is finite when minimal primes of $J$ are part of a specific regular system of parameters with a `Dropping Condition', see Corollary 1. This result is consistent with \cite{Ma}, for detail see Remark 2.%: Let $S$ be a %$d$-dimensional complete ramified regular local ring of mixed characteristic which is an Eisenstein extension of an unramified complete regular local ring $R$ sitting inside it. Let $J$ be an ideal of $S$. Assume that minimal primes of $J$ are the extension of ideals $I_t$ for $t=1,\ldots, n$ of $R$ such the quotient rings by $I_t S$'s are regular local rings. Then $\Ass H^i_{J} (S)$ is finite. These extended ideals $I_t S$'s are all cohomologically complete intersection \cite{HeSc}. 



%The proof of the main theorem uses %the properties of Eisenstein extension (see, for example \cite{Matsumura} for its reference) and  the Mayer-Vietoris spectral sequence (see \cite{AGZ}, \cite{ABZ} and \cite{Ly4}). %In ramified regular local ring for extended ideal (from the unramified one), set of associative primes of a local cohomology module is finite. Using the above result, for non extended ideal, we show example of local cohomology module which has finite set of associative primes and where associative primes also contains prime number $p$.
In Section 3 (The Main Result II), we prove Theorem 2.3 in more general sence, without the use of `Dropping Condition' and also the associated corollary. See Theorem 3.1 and its corollary. This corollary tells that for an ideal $J\subset S$ if its minimal primes are part of one regular system of parameters, then the set of associated primes of $H^i_J (S)$, is finite. This is not consistent with that of \cite{Ma}

Finally, in (\cite{Nu3}), it is shown that if the set of associated primes of a local cohomology module does not contain the prime number $p$, then the set of associated primes is finite. These examples of local cohomology modules (in this paper) include the cases whose associative primes also can contain prime number $p$, see Remark 5.

\section{the main result I}

%Let $(R;\m)$ be a local ring of mixed characteristic $p > 0$. We say $R$ is unramified if $p\notin {\m}^2$ and it is ramified if $p\in {\m}^2$. For normal local ring $(R,\m)$, consider the extension ring defined by $S=R[X]/f(X)$ where $f(X)=X^l+ a_1X^{l-1}+\ldots +a_l$ with $a_i\in \m$ for every $i=1,\ldots,l$ and $a_l\notin {\m}^2$. This ring $(S,\n)$ is local and it is defined as an Eisenstein extension of $R$ and $f(X)$ is known as an Eisenstein polynomial (see page 228-229 of \cite{Matsumura}). %More over from paragraph above Theorem 29.8 of \cite{Matsumura} we get that $\n= \m S+ xS \supset \m R[x] \supset {\m}^2$ and for $p\in \m$, $p\in {\n}^2$. Thus if $R$ is an unramified local ring then $S$ is a ramified local ring. 

%Here we note down the following important well known properties of Eisenstein extension:

%(1) An Eisenstein extension of regular local ring is regular local (see Theorem 29.8 (i) of \cite{Matsumura}) and in this context, we observe that an Eisenstein extensions of an unramified regular local ring is a ramified regular local ring.

%(2) Every ramified regular local ring is an Eisenstein extension of some unramified regular local ring (see Theorem 29.8 (ii) of \cite{Matsumura}. 

%(3) The Eisenstein extension mentioned in (1) and (2) are faithfully flat (apply Theorem 23.1 of \cite{Matsumura}). 


In proving our theorem, we need the Mayer-Vietoris spectral sequence, which arises from te Roos complex \cite{Roo} where in \cite{AGZ}, it is stated and proved only for defining ideal of an arrangement of linear subvarieties. In \cite{Ly4} it is extened for arbitrary commutative ring and module where \cite{ABZ}, the spectral sequnce has been constructed and properties are discussed very elaborately for more general categorical objects. Although the Mayer-Vietoris spectral sequences as observed in both \cite{AGZ} and \cite{Ly4}, differ in $E_0$ and $E_1$ page, but the coincide on $E_2$ page (see 4.2 Example \cite{ABZ} for detail).  

In this paper, we need the following theorem from \cite{Ly4} regarding the Mayer-Vietoris spectral sequence. 

\begin{thm}[Theorem on MVSS]\label{MV}
Let $A$ be an arbitrary commutative ring, let $I_1,\dots,I_n\subset A$ be ideals and
let $M$ be an $A$-module. There exists a spectral sequence
$$E_1^{-p,q}=\bigoplus_{i_0<\dots<i_p}H^q_{I_{i_0}+\dots+I_{i_p}}(M)
\Longrightarrow H^{q-p}_{I_1\cap\dots\cap I_n}(M).$$
\end{thm}

In our theorem we need the condition so that the above spectral sequence will be degenarate on $E_2$-page. In \cite{AGZ} in Theorem 1.2 it is shown the degenaracy in $E_2$-page for a defining ideal of an arrangement of linear subvarieties, when the underlying regular ring contains field. In \cite{ABZ} (see Lemma 4.15), degenaracy condition is analysed in a general setting and ideal which is cohomologically complete intersection (i.e. local cohomology module in the support of the ideal is non zero only at its height) used for the degenaracy. In \cite{Ly4} it has been mentioned that the spectral sequence will be degenarate on $E_2$ page if $A$ and $A/I_i$ are all regular (possibly arbitrary), with $M=A$ and gives reference of Theorem 1.2 of \cite{AGZ}.

To prove the main theorem, we also need minimal primes of the support ideal to be cohomologicaly complete intersection whose sums are also so. For an ideal $I$ in regular local ring $S$, if $S/I$ is regular then $I$ is cohomologicaly complete intersections (although converse is not true, consider $k[[x,y,z]]$ and the ideal $(x^2 +y^3 + z^5)$). 

%\begin{lemma}For a regular local ring $(A,\o)$ of dimension $n$, if $I, J$ are ideals generated by part of regular system of parameters, or equvalently $A/I, A/J$ are regular rings, such that $I+J$ is a prime, then $A/(I+J)$ is also a regular local ring. \end{lemma}
\begin{lemma}
Let $(A,\o)$ regular local ring.\\ 
(a) An ideal $K$ will be generated by regular system of parameters, or equivalently $A/K$ is regular if and only if $K\cap {\o}^2\subset \o K$.\\%$K/\o K$ is canonically isomorphic to $K/K\cap {\o}^2$. 
(b) If $I, J$ are ideals generated by part of regular system of parameters, or equivalently $A/I, A/J$ are regular rings and $A/(I+J)$ is also a regular then regular system of parameter of $I+J$ can get obtained by joining part of regular system of parameters of $I$ to that of $J$ and conversely if $I$ and $J$ are formed from a part of regular system of parameters then $I+J$ is generated by a regular system of parameters. 
\end{lemma}

\begin{proof}
(a) The asssertion will be true if and only if the minimal set of generators of $K$ is a regular system of parameters, since a regular system of parameters is a minimal set of generators for the ideal it generates. Now the minimala set of generators forms basis of the vector space $K/\o K$ over $A/\o$. So $K$ will be genetared by regular system of parameters if and only if the dimension of the vector space $(K+{\o}^2)/{\o}^2\cong K/K\cap {\o}^2$ is same to that of $K/\o K$ over $A/\o$. Since there is already a surjective canonical map from $K/\o K$ to $K/K\cap {\o}^2$, thus $K/\o K$ is canonically isomorphic to $K/K\cap {\o}^2$ if and only if $K\cap {\o}^2\subset \o K$. 

(b) Let $I=({u}_1,\ldots,{u}_k)A$ and $J=({v}_1,\ldots,{v}_l)A$ be generated by regular system of parameters and let $\{{u}_{i_1},\ldots,{u}_{i_a}, {v}_{j_1},\ldots,{v}_{j_b}\}$ be a set of minimal generatos of $I+J$. %Consider $(I+J)/{\o}(I+J)$ as vector space over $A/\o$ such that $\bar{u}_i$'s and $\bar{v}_j$'s spans it. Being minimal, both the sets of $\bar{u}_i$'s and $\bar{v}_j$'s are separately linearly independent, we can append part of $\bar{v}_j$'s to $\bar{u}_i$'s to get the basis of $(I+J)/{\o}(I+J)$. Suppose $I+J=({u}_1,\ldots,{u}_k, {v}_{t_1},\ldots,{v}_{t_b})A$ where $\{{u}_1,\ldots,{u}_k, {v}_{t_1},\ldots,{v}_{t_b}\}$ form the minimal set of generatos. 
From (a) of this lemma it follows that they will form the regular system of parameters and our assertion follows. Converse follows from the fact that joining two part of same regular system of parameters will be a part of regular system of parameters.
%Let $I=({u}_1,\ldots,{u}_k)A$ and $J=({v}_1,\ldots,{v}_l)A$ be generated by regular system of parameters and so $I+J=({u}_1,\ldots,{u}_k, {v}_1,\ldots,{v}_l)A$. Consider the ring $A/{\o}^2$, where sum of images of $I$ and $J$ forms a vector subspace $W\subset A/{\o}^2$ of dimension $t$ over $A/{\o}$. Suppose images $\bar{z}_1,\ldots,\bar{z}_t$ in $A/{\o}^2$ of the elements ${z}_1,\ldots,{z}_t$ from $\o$ form a basis of $W$. Thus $I+J+{\o}^2=({z}_1,\ldots,{z}_t)A +{\o}^2$. %Extend this basis to that of $\o/\o^2$ by the images $\bar{z}_1,\ldots,\bar{z}_t, \bar{z}_{t+1},\ldots, \bar{z}_n$ of the elements ${z}_1,\ldots,{z}_t,{z}_{t+1},\ldots, {z}_n$ from $\o$. Now, if any $z_i\in {\o}^2$ for $i=1,\dots,t$ dimension of subspace $W$ will be less than $t$. Infact there is another reason so that any $z_i$ can not be in ${\o}^2$, because they are all regular system of parameters. So all $({z}_1,\ldots,{z}_t)A\subset  I+J$. Similarly, $(I+J)\subset ({z}_1,\ldots,{z}_t)A +{\o}^2$ and since $u_i$'s and $v_j$'s ae regular system of parameters any of them can not be in ${\o}^2$. Thus $({z}_1,\ldots,{z}_t)A= I+J$. This proves the lemma. %Moreover, for any $k=t+1,\dots,n$, $z_{k}$ or any of its power can not be in $I+J$, otherwise being prime $I+J$ will contain $z_k$ and dimension of $W$ will be bigger than $t$. Thus $({z}_1,\ldots,{z}_t)A\subset  I+J +{\o}^2\subset ({z}_1,\ldots,{z}_t, z_{t+1})A$. Since ${z}_1,\ldots,{z}_n$ form regular system of parameters, both $({z}_1,\ldots,{z}_t)A$ and $({z}_1,\ldots,{z}_t, z_{t+1})A$ are primes along with $I+J$. Now $\height ({z}_1,\ldots,{z}_t)A= t$ while $\height ({z}_1,\ldots,{z}_t, z_{t+1})A= t+1$ and we have the following inclusions $({z}_1,\ldots,{z}_t)A\subset I+J \subset ({z}_1,\ldots,{z}_t, z_{t+1})A$ (where right inclusion is strict). This gives $({z}_1,\ldots,{z}_t)A=I+J$ and we conclude.
\end{proof}

Now we prove the main result of the paper.

\begin{thm}
Let $(S,\n)$ be a %$d$-dimensional 
complete ramified regular local ring of mixed characteristic and $J$ be an ideal of it. For $t=1,\ldots, n$, let $I_t$'s be the ideals  of $S$ and $J=\cap^n_{t=1} I_t$, such that all the ideals $\sum_{t\in F}I_t$ are prime and cohomologically complete intersection for every subsets $F\subset\{1,\dots,n\}$. Moreover we assume that dropping any $I_t$ from the $\sum_{t\in F}I_t$ gives proper containment. Then $\Ass H^i_{J} (S)$ is finite.
\end{thm}


\begin{proof}

The basic idea of the proof is as follows: Since for every subsets $F\subset\{1,\dots,n\}$, $\sum_{t\in F}I_t$ are all prime and cohomologically complete intersection, the Mayer-Vietoris spectral sequence will degenerate on $E_1$-page. This will imply that $H^i_{J} (S)$ has a finite filtration such that every subquotient module of the filtration is isomorphic to a local cohomology module whose set of associated primes is finite. Thus the set of associated primes of $H^i_{J} (S)$ is finite. 

We apply the above theorem on MVSS for the regular local ring $S$ and ideal $J$ with $M=S$. Here we start from $E_1$ page and define $E_1^{-p,q}$ according to \cite{Ly4}, infact previously, it has actually defined in a more general context in \cite{AGZ}. Set $J_{i_0\ldots i_p}=I_{i_0}+\dots+I_{i_p}$ and $J_{i_0\ldots\widehat{i_j}\ldots i_p}=I_{i_0}+\dots+I_{i_{j-1}}+I_{i_{j+1}}+\dots+I_{i_p}$:
$$E_1^{-p,q}=\bigoplus_{i_0<\dots<i_p}H^q_{J_{i_0\ldots i_p}}(S)
\Longrightarrow H^{q-p}_{I_1\cap\dots\cap I_n}(S)= H^{q-p}_{J}(S)$$

where differentials $E_1^{-p,q}\to E_1^{-p+1,q}$ on $E_1$-page are horizontal and they are defined according to \cite{Ly4} and in the following way: Consider their restriction to $H^q_{J_{i_0\ldots i_p}}(S)\subset E_1^{-p,q}$ where they take elements of it to some local cohomology module $H^{q}_{J_{i_0\ldots\widehat{i_j}\ldots,i_p}}(S)\subset E_2^{-p+1,q}$ and this map is induced by natural inclusion of $J_{i_0\ldots\widehat{i_j}\ldots i_p}\subset J_{i_0\ldots i_p}$. 

%At first, we observe the following: Let $P\subset Q$ be ideals of $S$ such that both $P$ and $Q$ are radical and cohomologically complete intersection. If $\height P= \height Q$ then $P=Q$. To see this let $\height P= \height Q= h$ and $H^h_P(S)$ and $H^h_Q(S)$. This is because any regular system of parameter of $P$ can be extend to $Q$. Indeed in the vector space $\n/{\n}^2$ over $S/{\n}$, basis of $(P+{\n}^2)/{\n}^2$ and $(Q+{\n}^2)/{\n}^2$ are same and their inverse images are the regular system of parameters where none of them can be inside ${\n}^2$.
%takes $x\in H^q_{I_{i_0} S+\dots+ I_{i_p} S}(S)\subset E_1^{-p,q}$ to$\oplus_{j=0}^{j=p}(-1)^jh^q_{J_{i_0\ldots i_p}, J_{i_0\ldots\widehat{i_j}\ldots i_p}}(x)$. Here $J_{i_0\ldots\widehat{i_j}\ldots i_p}=I_{i_0} S+\dots+I_{i_{j-1}} S+I_{i_{j+1}} S+\dots+I_{i_p} S$ and $J_{i_0\ldots i_p}=I_{i_0}S+\dots+I_{i_p}S$ while $h^q_{J_{i_0\ldots i_p}, J_{i_0\ldots\widehat{i_j}\ldots i_p}}:H^q_{J_{i_0\ldots i_p}}(S)\to H^q_{J_{i_0\ldots\widehat{i_j}\ldots i_p}}(S)$ is the natural map induced by the inclusion $J_{i_0\ldots\widehat{i_j}\ldots i_p}\subset J_{i_0\ldots i_p}$. %(hence $h^q_{J, J_j}(x)\in H^q_{J_j}(M)\subset E_1^{-p+1,q}$). 

If in $E_1$-page if $H^q_{J_{i_0\ldots i_p}}(S)$ is a non zero summand in $E_1^{-p,q}$ then $H^l_{J_{i_0\ldots\widehat{i_j}\ldots i_p}}(S)$ will be non zero for $l\leq q-1$. This is because we assume that dropping an index we have proper containment and since all $J_{i_0\ldots i_p}$ are primes, after dropping one index the height will be dropped by atleast one. Thus any nonzero local cohomology in the support of $J_{i_0\ldots\widehat{i_j}\ldots,i_p}$ appears along or below the diagonal of $H^q_{J_{i_0\ldots i_p}}(S)$.

As defined above, in $E_1$-page differentials are all zero. Thus on the $E_2$ page, we get $E_2^{-p,q}= E_1^{-p,q}=\bigoplus_{i_0<\dots<i_p}H^q_{J_{i_0\ldots i_p}}(S)$ along with the differentials $E_2^{-p,q}\to E_2^{-p+2,q-1}$. Consider a differential from $E_2^{-p,q}\to E_2^{-p+2,q-1}$ on the $E_2$-page and consider its restriction to $H^q_{J_{i_0\ldots i_p}}(S)\subset E_2^{-p,q}$ where they take elements of it to $E_2^{-p+2,q-1}=\bigoplus_{j_0<\dots<j_{p-2}}H^{q-1}_{J_{j_0\ldots j_{p-2}}}(S)$. If $H^{q-1}_{J_{j_0\ldots j_{p-2}}}(S)= 0$, then we have the zero map from $H^q_{J_{i_0\ldots i_p}}(S)$. Assume $H^{q-1}_{J_{j_0\ldots j_{p-2}}}(S)\neq 0$, this implies $\height J_{j_0\ldots j_{p-2}}= q-1$. Since $\height J_{i_0\ldots i_p}= q$, $J_{j_0\ldots j_{p-2}}$ can not contain $J_{i_0\ldots i_p}$ then according to Lemma 4.15 of \cite{ABZ}, there exists only the zero map from $H^q_{J_{i_0\ldots i_p}}(S)$ to $H^{q-1}_{J_{j_0\ldots j_{p-2}}}(S)$. So there is only zero differential in $E_2$-page.  
 

%Let $H^q_Q (R)$ be a nonzero direct summand of $E_1^{-p,q}$ and there is a map from here to $H^q_P (R)$ in $E_1^{-p+1,q}$ induced by the inclusion $P\subset Q$. Clearly $\height Q=q$. If $\height P=q$ then $Q=P$ and the map will be identity, otherwise for $\height P\neq q$, $H^q_P (R)=0$ and the map is a zero map. Now at each summand $H^q_{J_{i_0\ldots i_p}}(S)$ of $E_1^{-p,q}$ we map to $H^q_{J_{i_0\ldots\widehat{i_j}\ldots i_p}}(S)$ where $J_{i_0\ldots\widehat{i_j}\ldots i_p}\subset J_{i_0\ldots i_p}$. If $J_{i_0\ldots\widehat{i_j}\ldots i_p}= J_{i_0\ldots i_p}$ then the map is identity, which implies kernel of this differential is zero and this homology module can not appear in $E_2^{-p,q}$. Since modules on $E_2$-page are the homology modules of the differentials of $E_1$-page, on $E_2$-page, for any pair ideals $J_{i_0\ldots\widehat{i_j}\ldots i_p}$, $J_{i_0\ldots i_p}$ where  first one is obtained by deleting one index is can not be equal to the second one. If $J_{i_0\ldots\widehat{i_j}\ldots i_p}\subset J_{i_0\ldots i_p}$ properly then by above argument $\height J_{i_0\ldots\widehat{i_j}\ldots i_p}\leq q-1$. Thus $H^q_{J_{i_0\ldots\widehat{i_j}\ldots i_p}}(S)=0$ and the differential becomes a zero map on that particular summand of $E_1^{-p,q}$. Thus on the $E_2$ page, we get $E_2^{-p,q}\subset E_1^{-p,q}=\bigoplus_{i_0<\dots<i_p}H^q_{J_{i_0\ldots i_p}}(S)$ and in this $E_2$-page for any pair ideals $J_{i_0\ldots\widehat{i_j}\ldots i_p}$, $J_{i_0\ldots i_p}$ where  first one is obtained by deleting one index is contained in the second one properly. In fact, in $E_1$-page we left those pairs for which equality holds, since they give rise of identity maps between the corresponding local cohomologies, those local cohomologies do not appear in two places of $E_2$-page simaltaneously. So in $E_2$-page we fix the following notation: In $E_1$-page if $J_{i_0\ldots\widehat{i_j}\ldots i_p}= J_{i_0\ldots i_p}$, we call it as $J_{i_0\ldots\widehat{i_j}\ldots i_p}$ in $E_2$-page. 





% Here each containment $J_{i_0\ldots\widehat{i_j}\ldots,\widehat{i_k},\ldots,i_p}\subset J_{i_0\ldots\widehat{i_j}\ldots i_p}\subset J_{i_0\ldots i_p}$ is proper and hence $\height J_{i_0\ldots\widehat{i_j}\ldots,\widehat{i_k},\ldots,i_p}\leq q-2$ and hence $H^{q-1}_{J_{i_0\ldots\widehat{i_j}\ldots,\widehat{i_k},\ldots,i_p}}(S)=0$. So, there is no nonzero map between them. %to $\oplus_{j=0}^{j=p}(-1)^jh^q_{J_{i_0\ldots i_p}, J_{i_0\ldots\widehat{i_j}\ldots,\widehat{i_k},\ldots,i_p}}$ where $J_{i_0\ldots\widehat{i_j}\ldots,\widehat{i_k},\ldots,i_p}=I_{i_0} S+\dots+I_{i_{j-1}} S+I_{i_{j+1}} S+\dots+\dots+I_{i_{k-1}} S+I_{i_{k+1}} S+\dots+I_{i_p} S$. 
%To see this, suppose $J_{i_0\ldots i_p}= J_{i_0\ldots\widehat{i_j}\ldots,\widehat{i_k},\ldots,i_p}=J'$ (say), then atleast one of the local cohomology modules $H^q_{J'}(S)$ or $H^{q-1}_{J'}(S)$ will be a zero module. If $J_{i_0\ldots i_p}\supset J_{i_0\ldots\widehat{i_j}\ldots,\widehat{i_k},\ldots,i_p}$ where containment is proper, then by Lemma 4.5 of \cite{ABZ} (see also the Remark 4.16 of \cite{ABZ}) there will be no nonzero map. 
In fact for any positive integer $s$, inductively on $E_s$-th page, we have $E_{1}^{-p,q}= E_s^{-p,q}$ and by the same arguement as given to the $E_2$-page, the differential on the $E_s$-page is also zero. Thus $E_{1}^{-p,q}= E_{s+1}^{-p,q}$. Thus $E_{\infty}^{-p,q}= E_2^{-p,q}$. So for fixed $i=q-p$, $H^{i}_{J}(S)= H^{q-p}_{J}(S)= H^{q-p}_{I_1\cap\dots\cap I_n}(S)$ and we will get a filtration whose subquotient modules are isomorphic to $E_{\infty}^{-p,q}$ with varying $q$ and $p$ such that $i=q-p$. Since at each page of the spectral sequence we have only finitely many rows and columns, we have finite filtration and it is then sufficient to show that for each $(-p,q)$-pair, the set of asssociated primes of $E_{\infty}^{-p,q}$ is finite. From \cite{Mar} we know for every ideal $J_{j_0\ldots j_{p-2}}$ of $S$, $\Ass H^i_{J_{j_0\ldots j_{p-2}}}(S)$ is finite. So for each $(-p,q)$-pair $\Ass E_{\infty}^{-p,q}$ is finite. This implies $\Ass H^i_{J}(S)$ is contained in the subset of $\{J_{i_0\ldots i_p}:q-p=i\}$ and it is finite. This concludes the proof. 
\end{proof}

{\bf Definition}

In reference to the statement of Theorem 2.3, if we assume that dropping any $I_t$ from the $\sum_{t\in F}I_t$ gives proper containment, we define it as `Dropping Condition' satisfied by the minimal prime ideals.

%(It is to be noted that in the next manuscript "CCIinRR1.pdf" we have shown that for any ideal $I\subset S$ if $S/I$ is regular then set of associated primes of $H^i_I (S)$ is finite for all $i$ and we do not need the extended ideals from unramified regular local ring sitting inside it.)
{\bf Remark 1}\\

We can realize the above theorem in the following situation: For a ramified regular local ring $S$, let $x_1,\dots,x_d$ be a regular system of parameters such that $I_t$'s of the theorem are part of this regular system of parameters. From Lemma 2.2 it follows that the necessary and sufficient condition for the sum of any two such $I_t$ to be a part of regular system of parameters is to take all the $I_t$'s from one such regular system of parameters. Moreover we choose the $I_t$'s such that dropping any $I_t$ from the $\sum_{t\in F}I_t$ for every subsets $F\subset\{1,\dots,n\}$, gives proper containment. As for example for $n=3$, $I_1=(x_1,x_2)$, $I_2=(x_1,x_3)S$ and $I_3=(x_1,x_4)S$ taken from the regular sysytem of parameters $x_1,\ldots, x_d$, satisfy this condition. Thus we get the following corollary.

\begin{cor}
Consider a part of regular system of parameters $x_{i_1},\dots,x_{i_k}$ of a ramified regular local ring and consider ideals $I_t$ for $t=1,\ldots, n$ such that they are made out of parameters arbitrarily. Moreover we choose the $I_t$'s such that they satisfy the `Dropping Condition'. Then for the ideal $J=\cap^n_{i=1}I_t$ the set of associated primes of $H^i_{J}(S)$ is finite.
\end{cor}

{\bf Remark 2}\\

In \cite{Ma}, for $R = S/(2 + y^2_1 y^2_3 y_6 + y_1y_2y_3y_4y_5 + y^2_2 y^2_4 y_6)$ where $S = \mathbb{Z}_2[[x_1,\ldots, x_6, y_1,\ldots, y_6]] $ and for the ideal $J=IS +(y_1, y_2)S$ where $I\subset \mathbb{Z}_2[[x_1,\ldots, x_6]]$ is the square free monomial ideal corresponds to six vertex triangulation of $\mathbb{R}\mathbb{P}^2$, it is shown that local cohomology $H^6_J (R)$has infinitely many associated primes. Incidentally the set of minimal prime ideals of $J$ are $\{(\{x_i: i\notin H\}, y_1, y_2)\}$ (where $H$ varies over facets of the triangulation of $\mathbb{R}\mathbb{P}^2$), which part of regular system of parameters but the do not satisfy the `Dropping Condition' as defined above. 
%\begin{proof}From lemma 2.2, it follows that $I_t$'s satisfy the condition of the Theorem 2.3. So the result follows.
%\end{proof}


%%%%%%%%%%%%%%%%%%%%%%%%%%%%%%%%%%%%%%%%%%%%%%%%%%%%%%%%%%%%%%%%%%%%%%%%%%%%%%%%%%%%%%%%%%%%%%%%%%%%%%%%%%%%%%%%%%%%%%%%%%%%%%%%%%%%%%%%%%%%%%%%%%%%%%%%%%%%%%%%%%%%%%%%%%%%%%%%%%%%%%%%%%%%%%%%%%%%%%%%%%%%%%%%%%%%%%%%%%%%%%%%%%%%%%%%%%%%%%%%%%%%%%%%%%%%%%
\section{the main result II}

In this section, we prove Theorem 2.3 in more general sence, without the use of `Dropping Condition'. 

\begin{thm}
Let $(S,\n)$ be a %$d$-dimensional 
regular local ring and $J$ be an ideal of it. For $t=1,\ldots, n$, let $I_t$'s be the ideals  of $S$ and $J=\cap^n_{t=1} I_t$, such that all the ideals $\sum_{t\in F}I_t$ are prime and cohomologically complete intersection for every subsets $F\subset\{1,\dots,n\}$. Then $\Ass H^i_{J} (S)$ is finite.
\end{thm}


\begin{proof}

The basic idea of the proof is as follows: Since for every subsets $F\subset\{1,\dots,n\}$, $\sum_{t\in F}I_t$ are all prime and cohomologically complete intersection, the Mayer-Vietoris spectral sequence will degenerate on $E_2$-page. This will imply that $H^i_{J} (S)$ has a finite filtration such that every subquotient module of the filtration is isomorphic to a local cohomology module whose set of associated primes is finite. Thus the set of associated primes of $H^i_{J} (S)$ is finite. 

We apply the theorem on MVSS i.e. Theorem 2.1 above for the regular local ring $S$ and ideal $J$ with $M=S$. Here we start from $E_1$ page and define $E_1^{-p,q}$ according to \cite{Ly4}, infact previously, it has actually defined in a more general context in \cite{AGZ}. Set $J_{i_0\ldots i_p}=I_{i_0}+\dots+I_{i_p}$ and $J_{i_0\ldots\widehat{i_j}\ldots i_p}=I_{i_0}+\dots+I_{i_{j-1}}+I_{i_{j+1}}+\dots+I_{i_p}$ and from Theorem 2.1, we have the following convergence of the spectral sequence:
$$E_1^{-p,q}=\bigoplus_{i_0<\dots<i_p}H^q_{J_{i_0\ldots i_p}}(S)
\Longrightarrow H^{q-p}_{I_1\cap\dots\cap I_n}(S)= H^{q-p}_{J}(S)$$

Here differentials $E_1^{-p,q}\to E_1^{-p+1,q}$ on $E_1$-page are horizontal and they are defined according to \cite{Ly4} and in the following way: Consider their restriction to $H^q_{J_{i_0\ldots i_p}}(S)\subset E_1^{-p,q}$ where they take elements of it to some local cohomology module $H^{q}_{J_{i_0\ldots\widehat{i_j}\ldots,i_p}}(S)\subset E_2^{-p+1,q}$ and this map is induced by natural inclusion of $J_{i_0\ldots\widehat{i_j}\ldots i_p}\subset J_{i_0\ldots i_p}$. 

Let $H^q_Q (R)$ be a nonzero direct summand of $E_1^{-p,q}$ and there is a map from here to $H^q_P (R)$ in $E_1^{-p+1,q}$ induced by the inclusion $P\subset Q$. Clearly $\height Q=q$. If $\height P=q$ then $Q=P$ and the map will be identity, otherwise for $\height P\neq q$, $H^q_P (R)=0$ and the map is a zero map. Now at each summand $H^q_{J_{i_0\ldots i_p}}(S)$ of $E_1^{-p,q}$ we map to $H^q_{J_{i_0\ldots\widehat{i_j}\ldots i_p}}(S)$ where $J_{i_0\ldots\widehat{i_j}\ldots i_p}\subset J_{i_0\ldots i_p}$. If for every $i_j$, $J_{i_0\ldots\widehat{i_j}\ldots i_p}\subset J_{i_0\ldots i_p}$ properly, then $H^q_{J_{i_0\ldots\widehat{i_j}\ldots i_p}}=0$ and the restriction of the differential on $H^q_{J_{i_0\ldots i_p}}(S)$ will be a zero map. Now the differential map which ends on $E_2^{-p,q}$, the image of it in $H^q_{J_{i_0\ldots i_p}}(S)$, is either the whole $H^q_{J_{i_0\ldots i_p}}(S)$ or zero. Since the modules on $E_2$-page are the homology modules of the differentials of $E_1$-page and since homology commutes with the direct sum, in this case either $H^q_{J_{i_0\ldots i_p}}(S)$ appears or will not appear in $E_2^{-p,q}$ as a direct summand. On the other hand, if for any one $i_j$, $J_{i_0\ldots\widehat{i_j}\ldots i_p}= J_{i_0\ldots i_p}$ then the map is identity, which implies kernel of this differential restricted at $H^q_{J_{i_0\ldots i_p}}(S)$ is zero and this particular local cohomology summand module can not appear in $E_2^{-p,q}$. Thus on the $E_2$ page, we get $E_2^{-p,q}\subset E_1^{-p,q}=\bigoplus_{i_0<\dots<i_p}H^q_{J_{i_0\ldots i_p}}(S)$ and the differentials are from $E_2^{-p,q}\to E_2^{-p+2,q-1}$ on the $E_2$-page. Consider its restriction to $H^q_{J_{i_0\ldots i_p}}(S)\subset E_2^{-p,q}$ where they take elements of it to $E_2^{-p+2,q-1}\subset E_1^{-p+2,q-1}= \bigoplus_{j_0<\dots<j_{p-2}}H^{q-1}_{J_{j_0\ldots j_{p-2}}}(S)$. If $H^{q-1}_{J_{j_0\ldots j_{p-2}}}(S)= 0$, then we have the zero map from $H^q_{J_{i_0\ldots i_p}}(S)$. Assume $H^{q-1}_{J_{j_0\ldots j_{p-2}}}(S)\neq 0$, this implies $\height J_{j_0\ldots j_{p-2}}= q-1$. Since $\height J_{i_0\ldots i_p}= q$, $J_{j_0\ldots j_{p-2}}$ can not contain $J_{i_0\ldots i_p}$ then according to Lemma 4.15 of \cite{ABZ}, there exists only the zero map from $H^q_{J_{i_0\ldots i_p}}(S)$ to $H^{q-1}_{J_{j_0\ldots j_{p-2}}}(S)$. So there is only zero differential in $E_2$-page.  
 
In fact for any positive integer $s$, inductively on $E_s$-th page, we have $E_{2}^{-p,q}= E_s^{-p,q}$ and by the same arguement as given to the $E_2$-page, the differential on the $E_s$-page is also zero. Thus $E_{2}^{-p,q}= E_{s+1}^{-p,q}$. Thus $E_{\infty}^{-p,q}= E_2^{-p,q}$. So for fixed $i=q-p$, $H^{i}_{J}(S)= H^{q-p}_{J}(S)= H^{q-p}_{I_1\cap\dots\cap I_n}(S)$ and we will get a filtration whose subquotient modules are isomorphic to $E_{\infty}^{-p,q}$ with varying $q$ and $p$ such that $i=q-p$. Since at each page of the spectral sequence we have only finitely many rows and columns, we have finite filtration and it is then sufficient to show that for each $(-p,q)$-pair, the set of asssociated primes of $E_{\infty}^{-p,q}$ is finite. From \cite{Mar} we know for every ideal $J_{j_0\ldots j_{p-2}}$ of $S$, $\Ass H^i_{J_{j_0\ldots j_{p-2}}}(S)$ is finite. So for each $(-p,q)$-pair $\Ass E_{\infty}^{-p,q}$ is finite. This implies $\Ass H^i_{J}(S)$ is contained in the subset of $\{J_{i_0\ldots i_p}:q-p=i\}$ and it is finite. This concludes the proof. 
\end{proof}

%(It is to be noted that in the next manuscript "CCIinRR1.pdf" we have shown that for any ideal $I\subset S$ if $S/I$ is regular then set of associated primes of $H^i_I (S)$ is finite for all $i$ and we do not need the extended ideals from unramified regular local ring sitting inside it.)
{\bf Remark 3}\\

We can realize the above theorem in the following situation: For a regular local ring $S$, let $x_1,\dots,x_d$ be a regular system of parameters such that $I_t$'s of the theorem are part of this regular system of parameters. From Lemma 2.2 it follows that the necessary and sufficient condition for the sum of any two such $I_t$ to be a part of regular system of parameters is to take all the $I_t$'s from one such regular system of parameters. Thus we get the following corollary.

\begin{cor}
Let $S$ be a regular local ring and $J$ be its an ideal such that its minimal primes $I_t$ for $t=1,\ldots, n$ are part of regular system of parameters of a specific regular system of parameters $x_{1},\dots,x_{d}$ of $S$. Then the set of associated primes of $H^i_{J}(S)$ is finite.
\end{cor}

\begin{proof}
Since the ideal $J=\cap^n_{i=1}I_t$, from Remark 3 and Theorem 3.1 the result follows.
\end{proof}

{\bf Remark 4}\\

As discussed in Remark 2, this result is not consistent with \cite{Ma}.

{\bf Remark 5}\\

Consider the ramified regular local ring $S= {\mathbb{Z}}_p[[X_1,\dots, X_d]]/(X^2_1-p)$. Clearly in $S$, $x_1,\dots, x_d$ the images of $X_1,\dots, X_d$, form a regular system of parameters, such that $x^2_1=p$. If we take all the $I_t$'s which contains $x_1$, and satisfy the condition of the Corollary 1, then $J$ contains $x_1$ and hence $x^2_1= p$. From Corollary 1, we find that the set of associated primes of local cohomology modules $H^i_J (S)$ is finite and its associated primes contain $p$.
%%
%%%%%%%%%%%%%%%%%%%%%%%%%%%%%%%%%%%%%%%%%%%%%%%%%%%%%%%%%%%%%%%%%%%%%%%%%%%%%%%%%%%%%%%%%%%%%%%%%%%%%%%%%%%%%%%%%%%%%%%%%%%%%%%%%%%%%%%%%%%%%%%%%%%%%%%%%%%%%%%%%%%%%%%%%%%%%%%%%%%%%%%%%%%%%%%%%%%%%%%%%%%%%%%%%%%%%%%%%%%%%%%%%%%%%%%%%%%%%%%%%%%%%%%%%%%%%%%%%%%%%%%%%%%%%%%%%

\section{acknowledgement}

I would like to thank Professor Alberto F. Boix for pointing out errors in the earlier manuscript and providing the counter example that every cohomologically complete intersection ideal is not such an ideal that its quotient ring is regular local.

 \begin{thebibliography}{99}

\bibitem{AGZ} J. Alvarez Montaner, R. Garcia Lopez, and S. Zarzuela
Armengou, {\it Local cohomology, arrangements of subspaces and
monomial ideals}  Adv. Math.  {\bf 174}  (2003),  no. 1, 35--56. 

\bibitem{ABZ} J. Alvarez Montaner, Alberto F.Boix, and S. Zarzuela
Armengou, {\it On some Local cohomology spectral sequence}, International Mathematics Research Notices, Volume 2020, Issue 19, October 2020, Pages 6197?293

\bibitem{Boix} Alberto F.Boix, {\it Associated primes of local cohomology modules for five dimensional regular local rings}, DOI: 10.13140/RG.2.2.11073.52322

\bibitem{BBLSZ}
%\bibitem[BBLSZ]{BBLSZ}
B. Bhatt, M. Blickle, G. Lyubeznik, A. Singh and W. Zhang, {\it Local Cohomology Modules of a Smooth $\ZZ$-Algebra have Finitely Many Associated Primes}, Inventiones Mathematicae, 197 (2014) 509-519, arXiv: 1304.4692v2 [math.AC].
%\bibitem[BS98]{BroSharp}
%Markus~P. Brodmann and Rodney~Y. Sharp.
%\newblock {\em Local cohomology: an algebraic introduction with geometric
%  applications}, volume~60 of {\em Cambridge Studies in Advanced Mathematics}.
%\newblock Cambridge University Press, Cambridge, 1998.
%\bibitem[Bur]{Bur}
%L Burch.
%\newblock Codimension and andanalytic spread, 
%\newblock {\em Proc. Cambridge Philos. Soc.} 72, 1972, 369-373.

%\bibitem{HartshorneLC}
%\bibitem[Har67]{HartshorneLC}
%Robin Hartshorne.
%\newblock {\em Local cohomology}, volume 1961 of {\em A seminar given by A.
 % Grothendieck, Harvard University, Fall}.
%\newblock Springer-Verlag, Berlin, 1967.

%\bibitem{HartshorneCD}
%\bibitem[Har68]{HartshorneCD}
%Robin Hartshorne.
%\newblock Cohomological dimension of algebraic varieties.
%\newblock {\em Ann. of Math. (2)}, 88:403--450, 1968.

\bibitem{Hu}
%\bibitem[Hu1]{Hu}
C. Huneke, {\it Problems on local cohomology, in: Free resolutions in commutative algebra and algebraic
geometry (Sundance, Utah, 1990)}, 93-108, Res. Notes Math. 2, Jones and Bartlett, Boston, MA, 1992.

%\bibitem{H-L}
%\bibitem[HL90]{H-L}
%Craig Huneke and Gennady Lyubeznik.
%\newblock On the vanishing of local cohomology modules.
%\newblock {\em Invent. Math.}, 102(1):73--93, 1990.

\bibitem{HeSc}
%\bibitem[HeSc]{HeSc}
M. Hellus and P. Schenzel, {\it On cohomologically complete intersections}, Journal of Algebra, Volume 320 (10), 3733-3748, 2008. 

\bibitem{HS}
%\bibitem[HS]{HS} 
C. Huneke and R. Sharp, {\it Bass Numbers of Local Cohomology Modules}, Transaction of American Mathematical Society, Vol. 339, No. 2, 765-779, 1993.

%\bibitem{CohDim}
%\bibitem[HNBPW18]{CohDim}
%Daniel~J. Hern\'andez, Luis N\'u{\~n}ez-Betancourt, Juan~F. P\'erez, and
%  Emily~E. Witt.
%\newblock Cohomological dimension , Lyubeznik numbers, and connectedness in mixed characteristic
%\newblock {\em J. Algebra}, 514 (2018), 442-467
%\newblock {\em Preprint}, 2016.

\bibitem{Ka}
%\bibitem[Ka]{Ka}
M. Katzman, {\it An example of an infinite set of associated primes of a local cohomology module}, J. Algebra, 252, no. 1, 161?66, 2002.

%bibitem{Kwa}
%\bibitem{Kwa}[Kwa08] K. Kawasaki, {\it On finiteness properties of local cohomology modules over Cohen-Macaulay local ring} Illinois Journal of Mathematics Volume 52, Number 3, Fall 2008, Pages 727-744

\bibitem{Ly1}
%\bibitem[Ly1]{Ly1}
G. Lyubeznik, {\it Finiteness Properties of Local Cohomology modules (an application of D-modules to commutative algebra)}, Invent. Math. 113, 41-55, 1993.

\bibitem{Ly2}
%\bibitem[Ly2]{Ly2}
G. Lyubeznik, {\it Finiteness Properties of Local Cohomology Modules for Regular Local Rings of Mixed Characteristic: the Unramified Case}, Comm. Alg. 28, 5867-5882, 2000.

%\bibitem{Ly3}
%\bibitem[Ly3]{Ly3}G. Lyubeznik, {\it $F$-modules: applications to Local Cohomology and $D$-modules in Characteristic $p>0$}, J Reine Angew. Math. 491, 65-130, 1997.

\bibitem{Ly4}
%\bibitem{Ly4}[Lyu07] 
G. Lyubeznik, {\it On some local cohomology modules}, Advances in Mathematics Volume 213, Issue 2, 20 August 2007, Pages 621-643

\bibitem{Ma}
%\bibitem[Mar]{Mar}
L. Ma, \emph{Infinitely many associated primes of local cohomology modules of ramified regular local rings}, arXiv:  arXiv:2604.10964v2

\bibitem{Mar}
T. Marley, \emph{The associated primes of local cohomology modules over rings of small dimension}, Manuscripta Mathematica, Volume 104, pages 519-525 (2001).

%\bibitem{Matsumura}
%\bibitem[Mat80]{Matsumura} H. Matsumura \newblock {\it Commutative Ring Theory}, Cambridge University Press, 1990. 

\bibitem{Nu3}
%\bibitem[Nu3]{Nu3}
L. N$\acute{u}$$\tilde{n}$ez-Betancourt {\it On certain rings of differentiable type and finiteness properties of local cohomology}, Journal of Algebra, Volume 379, April 2013, Pages 1-10, arXiv: 1202.3597 [math.AC].
	
%\bibitem{Ogus}
%\bibitem[Ogu73]{Ogus}
%Arthur Ogus.
%\newblock Local cohomological dimension of algebraic varieties.
%\newblock {\em Ann. of Math. (2)}, 98:327--365, 1973.

%\bibitem{P-S}
%\bibitem[PS73]{P-S}
%C.~Peskine and L.~Szpiro.
%\newblock Dimension projective finie et cohomologie locale. {A}pplications \`a
%  la d\'emonstration de conjectures de {M}. {A}uslander, {H}. {B}ass et {A}.
%  {G}rothendieck.
%\newblock {\em Inst. Hautes \'Etudes Sci. Publ. Math.}, (42):47--119, 1973.

\bibitem{Roo} 
%\bibitem{Roo}[Roo]
J.-E. Roos, {\it Sur les foncteurs d$\acute{e}$riv$\acute{e}$s de lim. Applications}, C.R. Acad. Sci. Paris, 252:3702-3704, 1961. 

\bibitem{Si}
%\bibitem[Si]{Si}
A. K. Singh, {\it p-torsion elements in local cohomology modules}, Math. Res. Lett. 7, 165-176, 2000.

\bibitem{SS}
%\bibitem[SS]{SS}
A. K. Singh and I. Swanson, {\it Associated primes of local cohomology modules and of Frobenius powers}, Int. Math. Res. Not., no. 33, 1703-1733, 2004.
\end{thebibliography}


\end{document}
